\documentclass[11pt]{article}
\usepackage[a4paper,margin=25mm]{geometry}
\usepackage{amsmath,amssymb,amsthm,booktabs}
\usepackage[T1]{fontenc}
\usepackage{lmodern}
\usepackage[colorlinks=true,allcolors=blue]{hyperref}
\newtheorem{theorem}{Theorem}
\newtheorem{lemma}{Lemma}
\newcommand{\Q}{Q}
\title{Global optimality of every-jump resetting for a Cauchy flight:\\a computer-assisted survival comparison}
\author{Xiaojian Zeng\\\small PBC School of Finance, Tsinghua University, 100190\\\small\texttt{xiaojianzeng@foxmail.com}}
\date{}
\begin{document}
\maketitle
\begin{abstract}
For a discrete-time unit-Cauchy flight absorbed on landing above a barrier,
we prove that the two-step local condition for resetting before every jump
is sufficient for global optimality among constant Bernoulli resetting rates.
Specifically, if the reset-free path survival satisfies $Q_2(b)\ge F(b)^2$,
then every interior rate has a strictly larger mean first-passage time than
rate one. A finite set of outward-rounded killed-kernel inequalities controls
orders three through nine; the symmetric zero-barrier survival law controls
all higher orders. Separate integral enclosures exclude the antecedent at
all barriers above $1.809$. The result includes equality in the local
condition and arbitrarily small positive reset rates. It settles the Cauchy
specialization of the published global-optimality conjecture, without
asserting rate monotonicity or an extension to other stable indices.
\end{abstract}

\section{First passage and the endpoint question}
Resetting can improve first-passage performance, but an endpoint derivative
need not determine a global optimum. Radice, Cristadoro and Thapa
\cite{radice} derived a two-step criterion that excludes resetting before
every jump when its endpoint derivative has the wrong sign. Their stable-flight
discussion leaves the converse global conclusion conjectural. Here we prove
that converse for the Cauchy member of the family.

Discrete first passage under restart also has a general renewal description
\cite{bonomo}. Its first-passage probability generating function maps to the
survival generating function used here. We use that relation, rather than
claiming that the two descriptions necessarily have different reset timing.
The small-reset criterion in that work addresses the opposite endpoint;
it does not by itself establish the global every-jump implication below.
The zero-barrier symmetric survival law and the approximate Cauchy transition
barrier reported in Ref.~\cite{radice} are prior results, not contributions of
this calculation.

Let $X_0=0$ and let the independent increments have density and CDF
\begin{equation}
 f(x)=\frac{1}{\pi(1+x^2)},\qquad
 F(x)=\frac12+\frac{\arctan x}{\pi}.
\end{equation}
At each step an independent Bernoulli coin with probability $r$ resets the
position to zero \emph{before} the fresh increment:
\begin{equation}
 X_n=\begin{cases}\eta_n,&\text{with probability }r,\\
 X_{n-1}+\eta_n,&\text{with probability }1-r.
 \end{cases}
\end{equation}
For a finite barrier $b\ge0$, absorption occurs at
$\tau_b=\inf\{n\ge1:X_n>b\}$. Write $T_b(r)=\mathbb E\tau_b$.
Length is measured in the Cauchy scale and time in steps. The left half-line
is unbounded; no reflecting wall or finite-variance approximation is used.
Let $S_n=\sum_{j=1}^n\eta_j$ and define reset-free \emph{path} survival by
\begin{equation}
 Q_n(b)=\mathbb P(S_j\le b\text{ for every }1\le j\le n),\quad Q_0=1.
\end{equation}
In particular, $q=Q_1(b)=F(b)$ and
\begin{equation}\label{eq:q2}
 Q_2(b)=\int_{-\infty}^b f(x)F(b-x)\,dx.
\end{equation}
Endpoint $r=1$ gives independent trials, hence $T_b(1)=1/(1-q)$.

\begin{theorem}\label{thm:main}
For every finite $b\ge0$, in the process specified above,
\begin{equation}
 Q_2(b)\ge F(b)^2\quad\Longrightarrow\quad
 T_b(r)>\frac{1}{1-F(b)}\qquad(0<r<1).
\end{equation}
Thus $r=1$ is the unique minimum over constant rates $0<r\le1$ whenever
the antecedent holds, including its equality case.
\end{theorem}
The proof is computer assisted: the finite integral inequalities are
certified by integer interval arithmetic, whereas the infinite-order and
infinite-barrier extensions are analytic. It neither identifies every zero
of $Q_2-F^2$ nor proves monotonicity of $T_b(r)$.

\section{Renewal reduction and infinite orders}
Let $G_b(z)=\sum_{n\ge0}Q_n(b)z^n$ and $z=1-r$. For $0<r<1$ the renewal
mean is \cite{radice}
\begin{equation}\label{eq:mean}
 T_b(r)=\frac{zG_b(z)}{1-rG_b(z)}.
\end{equation}
For completeness, put $(Kv)(x)=\int_{-\infty}^b f(y-x)v(y)\,dy$.
The backward equation reads $u=1+zKu+r(Ku)(0)$.
With $h=(I-zK)^{-1}1$ one has $h(0)=G_b(z)$ and
$(Kh)(0)=(G_b(z)-1)/z$. Solving $u=ch$ for its scalar multiplier gives
Eq.~\eqref{eq:mean}. The killed reset transition has row mass at most
$1-r(1-q)<1$, so its mean exists and is bounded by $1/[r(1-q)]$.
Moreover $Q_n\le q$ for $n\ge1$ implies
\begin{equation}
 G_b(z)\le1+\frac{qz}{r},\qquad 1-rG_b(z)\ge z(1-q)>0.
\end{equation}
There is therefore no uncertain numerical denominator sign.

Direct subtraction in Eq.~\eqref{eq:mean} gives
\begin{equation}\label{eq:sign}
 \operatorname{sgn}[T_b(r)-T_b(1)]
 =\operatorname{sgn}[(1-zq)G_b(z)-1].
\end{equation}
It is enough that $Q_n(b)\ge q^n$ for every $n$, with at least one strict
inequality. Absolute convergence at $0<z<1$ then gives
$G_b(z)>1/(1-zq)$. This comparison avoids truncating the generating series
near $r=0$. Expanding its endpoint sign expression gives
$(1-zq)G_b(z)-1=(Q_2-q^2)z^2+O(z^3)$, explaining the local criterion.

Set $B=1809/1000$. The scalar enclosures described below prove
$F(B)<21/25$. For $0\le b\le B$, barrier inclusion and the continuous
symmetric survival law quoted in Ref.~\cite{radice} give
\begin{equation}
 Q_n(b)\ge Q_n(0)=a_n=4^{-n}\binom{2n}{n}.
\end{equation}
In exact arithmetic,
\begin{equation}
 a_{10}=\frac{46189}{262144}>\left(\frac{21}{25}\right)^{10},\qquad
 \frac{a_{n+1}}{a_n}=\frac{2n+1}{2n+2}\ge\frac{21}{22}>\frac{21}{25}
 \quad(n\ge10).
\end{equation}
Induction proves $Q_n(b)>q^n$ for all $n\ge10$. The zeroth and first
orders are identities; the second order is the theorem's hypothesis.
Only orders $3\le n\le9$ remain. As a convention check, $G_0(z)=(1-z)^{-1/2}$
in Eq.~\eqref{eq:mean} reproduces $T_0(r)=1+r^{-1/2}$.

\section{Uniform finite-order survival certificates}
For $x\le b$, let $u_n(x)$ be path survival for $n$ jumps starting at $x$.
Coupling the same increments proves that $u_n$ decreases in $x$, while
$Q_n(b)=u_n(0)$ increases in $b$. First-step conditioning gives
\begin{equation}\label{eq:kernel}
 u_{n+1}(x)=\int_{-\infty}^b f(y-x)u_n(y)\,dy,\qquad
 u_1(x)=F(b-x).
\end{equation}
Choose rational endpoints $-L=x_0<\cdots<x_m=b$, including zero.
If $\ell_{nk}\le u_n(x_k)$, then a lower node enclosure is
\begin{equation}\label{eq:lower}
 \ell_{n+1,j}=\sum_{k=1}^m
 \bigl[F(x_k-x_j)-F(x_{k-1}-x_j)\bigr]\ell_{nk},
\end{equation}
with every mass and sum rounded downward. The omitted integral below $-L$
is nonnegative. Discarding it weakens a lower bound; it does not introduce a
wall or presume that the omitted mass is small at every starting point.

A truncated approximation need not itself decrease. After each recurrence,
replace its values by $\max_{k\ge j}\ell_{nk}$. This remains a legal lower
bound at $x_j$, because true survival decreases. Thus the cell argument in
Eq.~\eqref{eq:lower} is valid at every order.

The certificate uses spacing $1/32$ from $-8$ to $b$, with geometric
negative cells from $-128$ to $-8$ and zero inserted as an endpoint.
For each closed barrier cell $[l,u]$, the checked inequalities are
\begin{equation}\label{eq:cell}
 Q_n(b)\ge Q_n(l)\ge\ell_n(l)>F(u)^n\ge F(b)^n,
 \qquad 3\le n\le9.
\end{equation}
The ten cells in Table~\ref{tab:cover} join exactly to cover $[0,B]$.
These are uniform event-inclusion inequalities, not interpolation of sampled
values. The table reports conservative decimal lower bounds on the smallest
margin in each cell; exact rational endpoints and margins accompany the code.
\begin{table}[ht]
\centering
\caption{Closed barrier cover for orders three through nine. Decimal margins are rounded downward; exact rational certificates are supplied.}\label{tab:cover}
\begin{tabular}{rrr}\toprule
$l$ & $u$ & Lower bound on minimum margin\\\midrule
0.00000000 & 0.45225000 & 0.05056029\\
0.45225000 & 0.67837500 & 0.06720767\\
0.67837500 & 0.90450000 & 0.04171904\\
0.90450000 & 1.13062500 & 0.02040270\\
1.13062500 & 1.35675000 & 0.00367261\\
1.35675000 & 1.46981250 & 0.01266770\\
1.46981250 & 1.58287500 & 0.00611999\\
1.58287500 & 1.69593750 & 0.00047858\\
1.69593750 & 1.75246875 & 0.00470379\\
1.75246875 & 1.80900000 & 0.00221400\\
\bottomrule\end{tabular}
\end{table}


CDFs are enclosed at integer scale $10^{30}$. The enclosure of $\pi$ uses
Machin's identity,
\begin{equation}
\pi=16\arctan(1/5)-4\arctan(1/239).
\end{equation}
Sign, reciprocal and addition identities reduce each rational arctangent
argument to absolute value at most $1/2$. Sixty alternating terms, together
with the next-term remainder, enclose the arctangent. Integer products,
divisions and positive recurrence sums round outwards or downwards as
appropriate. No floating-point transcendental evaluation contributes to a
certificate sign. Display decimals are not the decision arithmetic.

\section{Excluding every larger barrier}
Write $c_2(b)=Q_2(b)-F(b)^2$. We prove $c_2(b)<0$ for every $b\ge B$.
The three interval computations below are accompanied by exact rational
outputs and executable code.

\subsection{Initial sign and the interval to two}
To control the whole infinite integration domain in Eq.~\eqref{eq:q2}, set
$t=1/(1+b-x)$. Then
\begin{equation}\label{eq:compact}
 Q_2(b)=\int_0^1\frac{A(t)}{\pi D_b(t)}\,dt,\quad
 A(t)=F\left(\frac{1-t}{t}\right),\quad
 D_b(t)=t^2+[(b+1)t-1]^2.
\end{equation}
The endpoint value is $A(0)=1$. Composite Simpson quadrature with $N=8192$
and the explicit fourth-derivative bound in Appendix~\ref{app:error} gives,
rounding this display outward,
\begin{equation}
 -0.00000697205<c_2(B)<-0.00000694741.
\end{equation}

Differentiation of Eq.~\eqref{eq:q2} yields
$Q_2'(b)=f(b)/2+\int_{-\infty}^b f(x)f(b-x)\,dx$.
The full convolution is $2/[\pi(b^2+4)]$. The removed interval $[b,b+1]$
contributes at least $f(b+1)/4$, since its second Cauchy factor integrates
to $1/4$. Hence
\begin{equation}\label{eq:derivative}
 c_2'(b)\le\frac{2}{\pi(b^2+4)}
 -\frac{1}{4\pi[1+(b+1)^2]}
 -\frac{2F(b)-1/2}{\pi(1+b^2)}.
\end{equation}
For every cell $[l,u]$ of width $1/1000$ in $[B,2]$, evaluate the positive
term at $l$, the subtracted inverse quadratics at $u$, and the CDF factor
from below at $l$. All 191 resulting uniform upper bounds are negative.
The largest is less than $-0.00503$. Thus the initial sign persists to two.

\subsection{Crossing and returning between two and seven}
The marginal $S_2$ is Cauchy of scale two. Removing paths which crossed
at the first step and returned at the second gives the distinct identity
\begin{equation}\label{eq:return}
 Q_2(b)=F(b/2)-\int_0^\infty f(b+t)F(-t)\,dt.
\end{equation}
On a rational integration cell $[v,w]$, the removed integral is at least
$[F(b+w)-F(b+v)]F(-w)$. Summing these contributions with spacing $1/64$
to 16 and geometric cells to 1024 gives an upper enclosure of $Q_2$.
Discarding the remaining positive removed integral preserves that direction.
For a barrier cell $[l,u]$, use
\begin{equation}
 c_2(b)\le Q_2(u)-F(l)^2\qquad(l\le b\le u).
\end{equation}
Forty-two closed rational cells cover $[2,7]$, each with a strictly negative
upper bound; the largest is below $-0.0001724763$.

\subsection{The analytic infinite-barrier tail}
For $b\ge7$, put $x=1/b$, $t=\arctan x$ and $v=\arctan(2x)$.
Path inclusion implies $Q_2(b)\le F(b/2)$. On $0<x\le1/7$ there is no
arctangent branch ambiguity, and
\begin{align}
 F(b/2)-F(b)^2&=\frac{2t-v}{\pi}-\frac{t^2}{\pi^2},\\
 2t-v&=\arctan\frac{2x^3}{1+3x^2}\le2x^3,\qquad
 t\ge\frac{x}{1+x^2}.
\end{align}
The difference is strictly negative provided
$2\pi x(1+x^2)^2<1$. This increasing function is bounded at $x=1/7$
by $0.934609<1$. Consequently every barrier satisfying the theorem's
antecedent lies below $B$.

\section{Conclusion, interpretation and reproducibility}
On that remaining barrier interval, Eq.~\eqref{eq:cell} supplies orders
three through nine and the exact baseline induction supplies every order
at least ten. Together with the assumed second order, all survival
coefficients dominate $q^n$, strictly at order ten. Equation~\eqref{eq:sign}
therefore proves Theorem~\ref{thm:main} for every interior reset rate.

The physical consequence is a complete decision rule among constant
Bernoulli protocols: whenever the two-jump condition holds, resetting before
each jump is globally best. The equality case needs no numerical root
isolation. The proof controls the heavy left tail and the entire rate domain;
it makes no assumption of finite Cauchy moments. It does not settle the
other stable indices, uniqueness of a transition barrier, or optimization
over state-dependent or history-dependent protocols. A bounded current
literature check located no subsequent exact resolution of this Cauchy
implication; it does not guarantee priority.

The supplementary reproducibility instructions identify the exact integer
scripts, rational outputs and scalar checks. A separate reconstruction used
Horner interval arctangents, different scales and grids, and a legal
left-tail contribution; it passed all barrier cells and the initial sign.
That computational cross-check is not a substitute for human scrutiny of
the analytic enclosure arguments. The manuscript uses the accepted scope
only; the code contains more diagnostic outputs than the proof requires.

\paragraph{AI assistance and author review.}
Codex was used for research ideation, derivations, code generation, literature
checks, computational validation and manuscript drafting. Computational
checks include a separate AI-context reconstruction. No completed human
validation or authorship approval is asserted in this local draft. Before
submission, the named author must critically review the reasoning, code,
references and text, confirm the actual tool/model version and disclosure,
and take responsibility for the final work.

\appendix
\section{Explicit compact-quadrature error bound}\label{app:error}
Here is a reproducible bound for Eq.~\eqref{eq:compact}. Put
$a=1+(b+1)^2$, $d_1=\max\{2(b+1),2a-2(b+1)\}$ and $d_2=2a$.
On $[0,1]$, $D_b\ge1/a$, $|D_b'|\le d_1$ and $D_b''=d_2$.
Differentiating $D_b(1/D_b)=1$ gives valid bounds $B_k$ on the derivatives
of $1/D_b$ through order four:
\begin{equation}
 B_0=a,\qquad B_k=a\left[kd_1B_{k-1}+\binom{k}{2}d_2B_{k-2}\right],
\end{equation}
where the second term is omitted for $k=1$.
For $E=t^2+(1-t)^2$, $A'=-1/(\pi E)$, $E\ge1/2$,
$|E'|\le2$ and $E''=4$. Set
\begin{equation}
 C_0=2,\qquad C_k=2\left[2kC_{k-1}+4\binom{k}{2}C_{k-2}\right]
 \quad(1\le k\le3).
\end{equation}
With $\alpha_0=1$ and $\alpha_k=C_{k-1}/3$ for $1\le k\le4$,
$\pi>3$ and Leibniz's rule give
\begin{equation}
 \left|\left(\frac{A}{\pi D_b}\right)^{(4)}\right|
 \le M_4=\frac13\sum_{k=0}^4\binom4k\alpha_kB_{4-k}.
\end{equation}
All endpoint derivatives extend continuously. For even $N$ the composite
Simpson error is at most $M_4/(180N^4)$. The evaluated Simpson sum is itself
an outward interval; this rational error is added at both ends. No omitted
spatial-tail correction remains after the compact transformation.

\bibliographystyle{unsrt}
\bibliography{references}
\end{document}
